\title{Controlling Text-to-Image Diffusion by Orthogonal Finetuning}

\begin{abstract}
State-of-the-art text-to-image diffusion models demonstrate remarkable proficiency in synthesizing photorealistic visuals based on textual descriptions. However, a central unresolved issue is how to reliably steer or adjust these sophisticated models to handle a range of downstream applications. We address this gap by presenting a systematic finetuning approach—Orthogonal Finetuning~(OFT)—designed to tailor text-to-image diffusion models for various tasks beyond their original training. Distinct from prevalent approaches, OFT offers theoretical guarantees for maintaining hyperspherical energy, which encapsulates the relationships between neuron pairs on the unit hypersphere. Retaining this property proves essential for upholding the semantic generative capacity of text-to-image diffusion models. To further enhance the stability of the finetuning process, we introduce Constrained Orthogonal Finetuning (COFT), which incorporates an additional constraint on the hypersphere's radius. Our work focuses on two principal finetuning challenges: subject-driven generation, where the model must produce personalized images for a given subject based on a limited set of examples and a corresponding text prompt; and controllable generation, where the model is conditioned on extra control signals. Through empirical validation, we demonstrate that our OFT paradigm surpasses prior techniques in both quality of generation and rate of convergence.
\end{abstract}

\section{Introduction}

Recent advances in text-to-image diffusion models~\cite{saharia2022photorealistic,ramesh2022hierarchical,rombach2022high} have resulted in exceptional text-driven high-resolution image synthesis. Nonetheless, these systems often fall short in delivering precise and detailed control, as the information provided by text prompts can be ambiguous or incomplete. In this study, we concentrate on addressing two particular categories of text-to-image generation tasks:

\begin{itemize}[leftmargin=*,nosep]

    \item \textbf{Subject-driven generation}~\cite{ruiz2023dreambooth}: When only a few visual samples of a subject are available, the objective is to synthesize images depicting the same subject within new contexts, utilizing supplementary textual guidance.
    \item \textbf{Controllable generation}~\cite{zhang2023adding,mou2023t2i}: Given auxiliary control inputs (\emph{e.g.}, canny edges, segmentation maps), the task is to produce images that align with both these signals and a textual prompt.
\end{itemize}

The central challenge addressed by both tasks is how to efficiently adapt text-to-image diffusion models through finetuning, while retaining the generative strengths obtained during pretraining. We distill the key requirements for successful finetuning into two core aspects: (1) \emph{training efficiency}, encompassing reduced numbers of tunable parameters and fewer training iterations, and (2) \emph{generalizability preservation}, ensuring the model maintains high-quality and varied generation capabilities. Current finetuning approaches typically involve either incremental adjustment of neuron weights using a minimal learning rate (\emph{e.g.}, \cite{ruiz2023dreambooth}), or integrating a lightweight module that re-parameterizes certain weights (\emph{e.g.}, \cite{hulora2022,zhang2023adding}). Although straightforward, these strategies fall short of guaranteeing the retention of pretraining generative proficiency. Furthermore, there remains a gap in foundational insight regarding the systematic construction of finetuning procedures and the selection of suitable hyperparameters, such as the optimal number of training epochs. A primary challenge lies in the absence of a robust metric to assess the conservation of the original generative abilities post-finetuning. Present techniques often rely on the assumption that minimizing the Euclidean distance between finetuned and pretrained models equates to superior preservation of prior capabilities, which leads to the prevalent use of very small learning rates. While this notion occasionally proves effective, we contend that solely relying on Euclidean proximity does not adequately encapsulate semantic fidelity. Thus, incorporating a more structured measure to describe the relationship between the finetuned and pretrained models holds promise for enhancing both the maintenance of pretraining competence and the stability of the finetuning process.

Drawing on empirical findings that hyperspherical similarity effectively captures semantic relationships~\cite{liu2017deep,liu2018decoupled,chen2020angular}, we leverage hyperspherical energy~\cite{liu2018learning} to represent the pairwise relational patterns among neurons. Specifically, hyperspherical energy is computed as the aggregate of pairwise hyperspherical similarities (such as cosine similarity) among neurons within a single layer, which reflects the extent of uniformity among these neurons on the unit hypersphere~\cite{liu2021learning}. Our working hypothesis is that an optimally finetuned model should maintain a minimal alteration in hyperspherical energy relative to its pretrained counterpart. A straightforward strategy would involve adding a regularization term to preserve hyperspherical energy during finetuning, though this does not ensure effective minimization of the energy discrepancy. To address this, we exploit an important invariance property of hyperspherical energy: if every neuron undergoes the same orthogonal transformation, the pairwise hyperspherical similarity remains unchanged by construction. Building upon this invariant, we introduce Orthogonal Finetuning~(OFT), a technique for adapting large-scale text-to-image diffusion models to new tasks while keeping the hyperspherical energy constant. The method centers around learning a shared orthogonal transformation across neurons in each layer to preserve their mutual angular relationships. OFT can be interpreted as redefining the reference coordinate frame for neurons within a layer. Recognizing that reducing the Euclidean distance between the finetuned and pretrained models enhances the retention of pretraining benefits, we further present a constrained variant, termed Constrained Orthogonal Finetuning~(COFT), which restricts the adaptation such that the finetuned parameters reside on a hypersphere of fixed radius centered at the pretrained neurons.

The rationale behind employing orthogonal transformations in finetuning neurons is partially inspired by concepts from the 2D Fourier transform, a method that decomposes images into their magnitude and phase spectra. Of particular importance is the phase spectrum, representing angular relationships between the input and the basis, which retains the bulk of semantic content. Empirically, one can often reconstruct the semantic aspects of an image using its phase spectrum in conjunction with a randomized magnitude spectrum, highlighting that it is the orientation of neurons that fundamentally underlies semantic modification in generated images—a central challenge in both subject-driven and controllable generation. Yet, permitting excessive flexibility in altering neuron orientations can significantly impair the original pretrained generative capabilities. To address this, we advocate for constraining the freedom of transformation by maintaining the angular relationships between all pairs of neurons, guided by the premise that these inter-neuronal angles encapsulate core network knowledge. Building on this understanding, we propose applying a layer-wise, shared orthogonal transformation to the neurons in each layer, thereby ensuring that the hyperspherical energy remains constant.

Our approach also takes cues from orthogonal over-parameterized training~\cite{liu2021orthogonal}, in which classification networks are trained from scratch by imposing orthogonal transformations on a randomly initialized network. This stems from the observation that random initialization yields an expectedly low hyperspherical energy, and the method in~\cite{liu2021orthogonal} aims to sustain this low energy throughout training, as reduced hyperspherical energy has been linked to superior generalization in classification~\cite{liu2018learning,lin2020regularizing}. The results in~\cite{liu2021orthogonal} demonstrate that orthogonal transformations offer enough flexibility to train classification networks that generalize well. Our focus, however, is on refining text-to-image diffusion models to enhance control and generative efficacy on downstream tasks. We would like to underscore two key differences between OFT and the method in~\cite{liu2021orthogonal}. Firstly, while~\cite{liu2021orthogonal} is constructed to lower hyperspherical energy, OFT instead preserves the hyperspherical energy found in the pretrained model to prevent the degradation of intrinsic, learned structures during finetuning—indeed, reducing hyperspherical energy in diffusion model finetuning could undermine the existing semantic organization. Secondly, OFT prioritizes minimal deviation from the pretrained configuration, which motivates the introduction of a constrained formulation; in contrast,~\cite{liu2021orthogonal} enforces no such limits. We posit that a careful balance between flexibility and stability is vital for effective finetuning, and our OFT framework is deliberately crafted to strike this balance. Our main contributions are as follows:

\begin{itemize}[leftmargin=*,nosep]

    \item We introduce a new finetuning approach, Orthogonal Finetuning, aimed at enhancing the controllability of text-to-image diffusion models. To bolster stability during training, we also present a constrained version that restricts the angular change relative to the original pretrained weights.
    \item Distinct from prior finetuning strategies, OFT fine-tunes the model in a way that rigorously maintains the hyperspherical energy. Our experiments reveal this property to be critical for preserving the generative semantics encoded by the pretrained model.
    \item We employ OFT in both subject-driven and controllable generation tasks, undertaking extensive experimental analysis. Our results show notable gains compared to established baselines in terms of image quality, faster convergence, and stable finetuning dynamics. In addition, OFT demonstrates superior sample efficiency, achieving satisfactory performance with significantly fewer training samples and epochs.
    \item For the controllable generation task, we develop a novel control consistency metric to quantify the degree of controllability. This metric works by extracting the control signal from generated samples and comparing it against the original input control signal, thereby offering further evidence for the robustness of OFT in controllable image synthesis.
\end{itemize}

\section{Related Work}

\textbf{Text-to-image diffusion models}. The field of text-to-image synthesis has advanced rapidly~\cite{saharia2022photorealistic,ramesh2022hierarchical,rombach2022high,gu2022vector,nichol2022glide}, benefiting greatly from breakthroughs in both diffusion-based generative modeling~\cite{ho2020denoising,song2021score,song2021denoising,dhariwal2021diffusion} and joint vision-language representation learning~\cite{su2020vlbert,lu2019vilbert,radford2021learning,wang2021simvlm,li2022blip,li2023blip,alayrac2022flamingo,schuhmann2022laion}. Pioneering models such as GLIDE~\cite{nichol2022glide} and Imagen~\cite{saharia2022photorealistic} operate diffusion in pixel space. In these, GLIDE combines training of the text encoder and diffusion model with paired data, while Imagen relies on a frozen, pretrained language model. In contrast, latent diffusion approaches such as Stable Diffusion~\cite{rombach2022high} and DALL-E2~\cite{ramesh2022hierarchical} operate on latent representations: Stable Diffusion leverages VQ-GAN~\cite{esser2021taming} for a visual tokenization, whereas DALL-E2 uses CLIP~\cite{radford2021learning} to align image and text embeddings. Besides diffusion, text-to-image research includes generative adversarial networks~\cite{reed2016generative,zhang2017stackgan,xu2018attngan,li2019controllable} and autoregressive models~\cite{ramesh2021zero,ding2021cogview,wu2022nuwa,yuscaling}. OFT itself is a generic finetuning scheme and applicable to any text-to-image diffusion architecture.

\textbf{Subject-driven generation}. Techniques aimed at minimizing subject alteration, such as those in~\cite{avrahami2022blended,nichol2022glide}, introduce user-provided masks as supplemental conditioning. Alternative inversion-based methods~\cite{choi2021ilvr,dhariwal2021diffusion,ramesh2022hierarchical,gal2022image} enable modifying image context while maintaining the main subject. The framework in~\cite{hertz2022prompt} enables local and global edits without explicit masks. However, these methods generally fall short in preserving intricate identity features. Pivotal Tuning~\cite{roich2022pivotal} circumvents this by finetuning a generator near an initial inverted code and adding identity-preserving regularization. In a related vein, \cite{nitzan2022mystyle} constructs a tailored generative prior using an individual's facial images. Although \cite{casanova2021instance} offers instance variability, it may lose subject-specific information. DreamBooth~\cite{ruiz2023dreambooth} instead adapts a diffusion model using a personalized token and a handful of subject images, optimizing with both a reconstruction loss and a class-preserving prior loss. In our approach, we retain the DreamBooth methodology but replace standard finetuning with orthogonal transformations via OFT, moving beyond simple small-step updates.

\textbf{Controllable generation}. Image-to-image translation represents a subtype of controllable generation, wherein prior research predominantly leverages conditional generative adversarial networks~\cite{isola2017image,zhu2017toward,wang2018high,park2019semantic,choi2018stargan}. Recent advancements have seen diffusion models applied to the image-to-image translation task~\cite{saharia2022palette,voynov2022sketch,wang2022pretraining}. In this domain, ControlNet~\cite{zhang2023adding} introduced an approach that steers a pretrained diffusion model through finetuning with supplementary control signals, demonstrating robust controllable generation capabilities. Similarly, T2I-Adapter~\cite{mou2023t2i}, a contemporaneous method, refines a pretrained diffusion model to augment the controllability of synthesized outputs. Adhering to the experimental protocols outlined in \cite{zhang2023adding,mou2023t2i}, we deploy OFT to finetune diffusion models pretrained on large datasets. This approach not only enhances control over the generative process but does so using less training data and fewer finetuned parameters. Importantly, OFT maintains inference-time efficiency without incurring extra computational cost.

\textbf{Model finetuning}. The practice of adapting large-scale pretrained models to downstream tasks via finetuning has become increasingly widespread~\cite{devlin2018bert,brown2020language,he2022masked}. Various adaptation techniques, such as those introduced in \cite{hulora2022,houlsby2019parameter,pfeiffer2020adapterfusion}, have received considerable attention, particularly in natural language processing. Of special relevance is LoRA~\cite{hulora2022}, which models finetuning updates as low-rank modifications to weights. Distinctly, OFT employs a multiplicative, orthogonal transformation—shared across layers—to modify neuron weights. This formulation ensures the preservation of pairwise angles between neurons within each layer, leading to greater training stability.

\section{Orthogonal Finetuning}

\subsection{Why Does Orthogonal Transformation Make Sense?}

To motivate the use of orthogonal transformation in the context of finetuning text-to-image diffusion models, we decompose the rationale into two central considerations: (1) the significance of refining neuron angles (\emph{i.e.}, their directional properties), and (2) the suitability of orthogonal transformation as a mechanism for this refinement.

Addressing the first question, we take cues from the empirical findings in \cite{liu2018decoupled,chen2020angular}, which suggest that differences in angular features provide a strong characterization of semantic gaps. According to SphereNet~\cite{liu2017deep}, constraining all neurons of a neural network to lie on the surface of a unit hypersphere during training retains the network’s capacity and can even improve generalization, suggesting that neuronal directionality encapsulates essential data information. To underscore the significance of neuronal angles, we set up a simple experiment, illustrated in Figure~\ref{fig:toy_ae}, where a standard convolutional autoencoder is trained on flower image datasets. During the training phase, we use the conventional inner product for feature map extraction ($z$ represents the output element from convolution with kernel $\bm{w}$ applied to input $\bm{x}$ within the receptive field). At test time, we evaluate three feature generation strategies: (a) retaining the inner product from training, (b) considering only magnitude, and (c) utilizing only the angular component. The outcomes displayed in Figure~\ref{fig:toy_ae} reveal that recovering images using angular neuron information is nearly perfect, whereas magnitude alone proves uninformative. Notably, cosine-based angular activations (c) are not included during training, which relies exclusively on inner products. These observations indicate that neuron angles—i.e., directions—predominantly encode the semantic details of the inputs. Thus, when aiming to alter semantic content in images, adjusting neuron directions may offer the most effective route.

For the second question, one straightforward approach to refining neuron direction involves collectively rotating or reflecting all neurons within the same layer—operations that inherently utilize orthogonal transformations. While it is possible to consider more general angular modifications, such as rotating individual neurons by unique angles, our findings indicate that orthogonal transformations achieve a practical balance between adaptability and structural constraint. Additionally, as shown in \cite{liu2021orthogonal}, such transformations are sufficiently expressive for neural network learning tasks. To empirically support this, we conduct an experiment to highlight the regularization effect introduced by enforcing orthogonality. Using the controllable image generation setup from ControlNet~\cite{zhang2023adding}, with results shown in Figure~\ref{fig:ortho}, we observe that applying standard OFT (Orthogonal Finetuning) yields stable performance and precise controllability upon convergence (epoch 20). In contrast, removing the orthogonality requirement from OFT leads to the inability to produce realistic images or achieve control. This experiment affirms the critical role of the orthogonality constraint within OFT.

\subsection{General Framework}

In standard finetuning approaches, gradient descent is commonly employed to adjust the model parameters (or some of its layers), typically with a relatively low learning rate. This conservative update mechanism implicitly limits how much the parameters deviate from those of the original pretrained model; the finetuning procedure is thus approximately equivalent to minimizing $\thickmuskip=2mu \medmuskip=2mu \| \bm{M} - \bm{M}^0\|$, where $\bm{M}$ denotes the updated weights and $\bm{M}^0$ stands for the weights prior to finetuning. While intuitively reasonable, this implicit regularization might not be stringent enough when dealing with large models. To better regulate the adaptation, LoRA proposes a supplementary low-rank restriction on the difference between weights, specifically, $\thickmuskip=2mu \medmuskip=2mu \text{rank}(\bm{M} - \bm{M}^0)=r'$ with $r'$ chosen to be a small value. In contrast, OFT implements a restriction focused on preserving the pair-wise similarity among neurons by enforcing $\thickmuskip=2mu \medmuskip=2mu \| \text{HE}(\bm{M})-\text{HE}(\bm{M}^0) \|=0$, where $\text{HE}(\cdot)$ measures the hyperspherical energy of the respective weight matrices. For illustration, take a fully connected layer parameterized by $\thickmuskip=2mu \medmuskip=2mu \bm{W}=\{\bm{w}_1,\cdots,\bm{w}_n\}\in\mathbb{R}^{d\times n}$, with each neuron represented by $\bm{w}_i\in\mathbb{R}^{d}$ (and $\bm{W}^0$ representing its pretrained counterpart). The output of this layer, given input $\thickmuskip=2mu \medmuskip=2mu \bm{x}\in\mathbb{R}^d$, is $\thickmuskip=2mu \medmuskip=2mu \bm{z}=\bm{W}^\top\bm{x}$ where $\thickmuskip=2mu \medmuskip=2mu \bm{z}\in\mathbb{R}^n$. The core principle of OFT is then to minimize the disparity in hyperspherical energy between the adapted and original weight matrices:

\begin{equation}\label{eq:oft_principle}
\footnotesize
\min_{\bm{W}} \left\| \text{HE}(\bm{W})-\text{HE}(\bm{W}^0)\right\|~~\Leftrightarrow~~\min_{\bm{W}} \bigg{\|} \sum_{i\neq j} \|\hat{\bm{w}}_i-\hat{\bm{w}}_j\|^{-1}-\sum_{i\neq j} \|\hat{\bm{w}}^0_i-\hat{\bm{w}}^0_j\|^{-1}\bigg{\|}
\end{equation}

where the normalized neuron is defined as $\thickmuskip=2mu \medmuskip=2mu \hat{\bm{w}}_i:=\bm{w}_i/\|\bm{w}_i\|$, and the hyperspherical energy of $\bm{W}$ is given by $\thickmuskip=2mu \medmuskip=2mu \text{HE}(\bm{W}):= \sum_{i\neq j} \|\hat{\bm{w}}_i-\hat{\bm{w}}_j\|^{-1}$. The objective in Eq.~\eqref{eq:oft_principle} attains its minimum value of zero when $\bm{W}$ and $\bm{W}^0$ are related by an orthogonal transformation, i.e., $\thickmuskip=2mu \medmuskip=2mu \bm{W}=\bm{R}\bm{W}^0$ for some orthogonal matrix $\thickmuskip=2mu \medmuskip=2mu \bm{R}\in\mathbb{R}^{d\times d}$ (with determinant $1$ or $-1$ indicating whether it is a rotation or reflection). This constitutes the main motivation behind OFT: finetuning is achieved by learning orthogonal matrices (shared across layers) to rotate or reflect the neurons within each layer. In summary, the OFT approach focuses

\begin{equation}\label{eq:oft_general}
    \footnotesize
    \bm{z}=\bm{W}^\top\bm{x}=(\bm{R}\cdot\bm{W}^0)^\top\bm{x},~~~\text{s.t.}~\bm{R}^\top\bm{R}=\bm{R}\bm{R}^\top=\bm{I}
\end{equation}

Here, $\bm{W}$ represents the weight matrix obtained after OFT finetuning, while $\bm{I}$ is the identity matrix. A visualization of OFT is provided in Figure~\ref{fig:block}. Drawing a parallel to the zero initialization approach used in LoRA, it is necessary to configure OFT such that the finetuning process originates precisely from the pretrained parameters $\bm{W}^0$. To ensure this behavior, we initialize the orthogonal matrix $\bm{R}$ as the identity matrix, allowing the finetuned model to begin with the original pretrained weights. To maintain the orthogonal constraint on $\bm{R}$, various differentiable orthogonalization techniques, as outlined in \cite{liu2021orthogonal,lezcano2019cheap}, can be leveraged. We will elaborate on methods to uphold orthogonality efficiently from a computational perspective.

\subsection{Efficient Orthogonal Parameterization}\label{sect:eff_ortho}

While conventional orthogonalization procedures like the Gram-Schmidt process are differentiable, they are often computationally burdensome in real-world settings~\cite{liu2021orthogonal}. To enhance efficiency, we utilize Cayley parameterization for constructing the orthogonal matrix. Specifically, this approach forms $\bm{R}$ as $\thickmuskip=2mu \medmuskip=2mu \bm{R}=(\bm{I}+\bm{Q})(I-\bm{Q})^{-1}$, where $\bm{Q}$ denotes a skew-symmetric matrix satisfying $\thickmuskip=2mu \medmuskip=2mu \bm{Q}=-\bm{Q}^\top$. Although Cayley parameterization is limited to orthogonal matrices with determinant $1$, thus restricting the output to the special orthogonal group, our observations indicate that this does not compromise empirical performance. Nevertheless, parameterizing $\bm{R}$ with Cayley transform can still result in significant parameter overhead, particularly for large $d$. To tackle this, we recommend structuring $\bm{R}$ as a block-diagonal matrix with $r$ blocks, which leads to the following expression:

\begin{equation}
    \footnotesize
    \bm{R}=\text{diag}(\bm{R}_1,\bm{R}_2,\cdots,\bm{R}_r)=\begin{bmatrix}
    \bm{R}_1\in \text{O}(\frac{d}{r}) &  &\\
    &\ddots & \\
    & & \bm{R}_r\in \text{O}(\frac{d}{r})
    \end{bmatrix}\in \text{O}(d)
\end{equation}

Here, $\text{O}(d)$ represents the orthogonal group of degree $d$. Both $\thickmuskip=2mu \medmuskip=2mu \bm{R}\in\mathbb{R}^{d\times d}$ and each block matrix $\thickmuskip=2mu \medmuskip=2mu \bm{R}_i\in\mathbb{R}^{d/r\times d/r},~\forall i$ are required to be orthogonal. If $\thickmuskip=2mu \medmuskip=2mu r=1$, then the block-diagonal orthogonal structure collapses to a conventional full orthogonal matrix without any block constraint. The parameter count for a general $d\times d$ orthogonal matrix is $\thickmuskip=2mu \medmuskip=2mu d(d-1)/2$, entailing a complexity of $\mathcal{O}(d^2)$. In contrast, for an $r$-block diagonal orthogonal matrix, there are $\thickmuskip=2mu \medmuskip=2mu d(d/r-1)/2$ parameters, which yields a reduced complexity of $\mathcal{O}(d^2/r)$. Additionally, if all blocks are shared—specifically, by letting $\thickmuskip=2mu \medmuskip=2mu\bm{R}_i=\bm{R}_j$ for all $i\neq j$—the number of parameters decreases further to an order of $\mathcal{O}(d^2/r^2)$. It is important to emphasize that all of these approaches maintain the orthogonality of $\bm{R}$, ensuring no compromise in the preservation of hyperspherical energy.

We now compare OFT and LoRA concerning parameter efficiency. LoRA with rank $r'$ employs $\thickmuskip=2mu \medmuskip=2mu r'(d+n)$ trainable parameters. Setting both $r$ and $r'$ proportional to $d$—such as $\thickmuskip=2mu \medmuskip=2mu r=r'=\alpha d$ for some $0<\alpha\leq 1$—yields LoRA’s complexity as $\mathcal{O}(d^2+dn)$, while OFT’s complexity reduces to $\mathcal{O}(d)$. To illustrate, consider a weight matrix of dimensions $\thickmuskip=2mu \medmuskip=2mu 128\times 128$; with $\thickmuskip=2mu \medmuskip=2mu r'=8$, LoRA requires $2{,}048$ tunable parameters, while OFT, using $r=8$ and no block sharing, involves only $960$ parameters.

\subsection{Constrained Orthogonal Finetuning}

Original OFT’s representational capacity can be further regulated by constraining the adaptation to remain close to the pretrained model. In Constrained Orthogonal Finetuning (COFT), the forward computation is formulated as

\begin{equation}\label{eq:coft_general}
    \footnotesize
    \bm{z}=\bm{W}^\top\bm{x}=(\bm{R}\cdot\bm{W}^0)^\top\bm{x},~~~\text{s.t.}~\bm{R}^\top\bm{R}=\bm{R}\bm{R}^\top=\bm{I},~\norm{\bm{R}-\bm{I}}\leq \epsilon
\end{equation}

where $\bm{R}$ must remain both orthogonal and within an $\epsilon$-distance of the identity matrix. The orthogonality is preserved through Cayley parameterization (see Section~\ref{sect:eff_ortho}). Imposing the additional $\epsilon$-proximity constraint on the Cayley-parameterized matrix is, however, challenging. To explore the characteristics of the Cayley transform, we use the Neumann series to approximate $\thickmuskip=2mu \medmuskip=2mu \bm{R}=(\bm{I}+\bm{Q})(I-\bm{Q})^{-1}$ by expanding as $\thickmuskip=2mu \medmuskip=2mu \bm{R}\approx\bm{I}+2\bm{Q}+\mathcal{O}(\bm{Q}^2)$, under the Neumann series assumption.

Since the series converges in the operator norm, the constraint $\thickmuskip=2mu \medmuskip=2mu\|\bm{R}-\bm{I}\|\leq\epsilon$ can be equivalently reformulated within the Cayley transform as the condition $\thickmuskip=2mu \medmuskip=2mu \|\bm{Q}-\bm{0}\|\leq\epsilon'$, where $\bm{0}$ represents the zero matrix and $\epsilon'$ is a distinct error tolerance parameter from $\epsilon$. Enforcing this alternative constraint on $\bm{Q}$ is straightforward using projected gradient descent. For the orthogonal matrix $\bm{R}$ to be initialized as the identity, we simply initialize $\bm{Q}$ to the zero matrix. The COFT method effectively introduces two explicit regularizations: limiting the change in hyperspherical energy and explicitly constraining the deviation from the original pretrained model. Although traditional finetuning strategies often enforce the latter constraint only implicitly, COFT makes this regularization direct and explicit. Our empirical analysis suggests that while OFT performs robustly, the additional explicit deviation constraint in COFT further enhances the stability of finetuning. Figure~\ref{fig:eps_coft} visualizes the impact of varying $\epsilon$ on COFT's performance, showing that $\epsilon$ determines the adaptability of finetuning—higher values make the COFT-adapted model behave similarly to the OFT-finetuned version, whereas lower values push the model closer to the behavior of the original text-to-image diffusion network.

\subsection{Re-scaled Orthogonal Finetuning}

We extend the baseline OFT method by introducing a learnable scaling factor for each neuron. The primary motivation for this is that scaling neurons does not influence the hyperspherical energy, as normalization counters any magnitude change. Formally, we define the forward pass as: $\thickmuskip=2mu \medmuskip=2mu\bm{z}=(\bm{R}\bm{W}^0\bm{D})^\top\bm{x}$\footnote{\textbf{Errata}: The NeurIPS camera-ready submission incorrectly listed the re-scaled OFT forward pass as $\thickmuskip=2mu \medmuskip=2mu\bm{z}=(\bm{D}\bm{R}\bm{W}^0)^\top\bm{x}$. The implementation, however, is correct, so Appendix~\ref{app:rescaled} reports valid results.}, where $\thickmuskip=2mu \medmuskip=2mu \bm{D}=\text{diag}(s_1,\cdots,s_n)\in\mathbb{R}^{n\times n}$ denotes a diagonal matrix with positive, learnable entries $s_1,\cdots,s_n$. While original OFT only adapts $\bm{R}$ as described in Eq.~\eqref{eq:oft_general}, this extension allows both the diagonal scaling $\bm{D}$ and the orthogonal matrix $\bm{R}$ to be jointly optimized. The added flexibility brought by this re-scaling comes with minimal increase in parameter count. In our main experiments, we focus on the standard OFT formulation to isolate the effect of the orthogonal transformation, but observe that the re-scaled variant consistently outperforms it (refer to Appendix~\ref{app:rescaled}).

\section{Intriguing Insights and Discussions}

\textbf{OFT is universally applicable across architectures.} OFT is not limited to any specific model structure and can theoretically be applied to any neural network. For example, while LoRA typically adapts the attention weights in Transformers~\cite{hulora2022}, for an equitable comparison, we restrict OFT’s application to the same set of weights in our benchmarks. Beyond fully connected layers, OFT’s structure is particularly advantageous for convolutional layers—unlike LoRA—since its block-diagonal orthogonal matrices can be naturally interpreted in the convolutional setting. If the number of blocks equals the number of input channels, each block operates on a distinct neuron channel, reminiscent of depth-wise convolutions~\cite{chollet2017xception}. If blocks are shared across channels, OFT applies a common orthogonal transformation to all neuron channels. Preliminary investigations of OFT applied to convolution layers are discussed in Appendix~\ref{app:convolution}.

\textbf{Connection to LoRA}. LoRA introduces a low-rank matrix to ensure that the original information in the pretrained weight matrix is largely retained. In contrast, OFT enforces the transformation applied to the pretrained weights to be orthogonal (and thus full-rank), safeguarding the pretrained information from being degraded by the transformation. The forward computation in OFT can be reformulated as $\thickmuskip=2mu \medmuskip=2mu \bm{z}=(\bm{R}\bm{W}^0)^\top\bm{x}=(\bm{W}^0+(\bm{R}-\bm{I})\bm{W}^0)^\top\bm{x}$, where $\thickmuskip=2mu \medmuskip=2mu (\bm{R}-\bm{I})\bm{W}^0$ plays a role similar to LoRA's low-rank adaptation. Given that $\bm{W}^0$ is usually full-rank, OFT also amounts to a low-rank modification of the weights whenever $\thickmuskip=2mu \medmuskip=2mu \bm{R}-\bm{I}$ itself is low-rank. Comparable to LoRA’s use of a rank hyperparameter $r'$, OFT employs a block-diagonal size parameter $r$ to manage model efficiency. Importantly, OFT and LoRA approach parameter efficiency via fundamentally different mechanisms: LoRA leverages low-rank structures to reduce trainable parameters, while OFT achieves efficiency by capitalizing on block-diagonal (sparse) orthogonal transformations.

\textbf{Why OFT converges faster}. Evidence from Figure~\ref{fig:toy_ae} shows that the primary mechanism for updating semantic information involves altering neuron orientations, which is precisely the modification that OFT targets. Furthermore, the OFT framework can be interpreted as performing finetuning of neural representations constrained on a smooth hyperspherical manifold, resulting in a more favorable optimization landscape. This benefit is also empirically demonstrated in \cite{liu2021orthogonal}.

\textbf{Why not minimize hyperspherical energy}. Our approach departs from \cite{liu2021orthogonal} in that we do not set out to minimize hyperspherical energy. In classification tasks, it is preferable for neurons to remain non-redundant; minimizing hyperspherical energy would lead to a uniform distribution of neurons on the hypersphere, which could undermine useful structure from pretraining. Such a configuration is thus inappropriate for finetuning, as it risks discarding valuable pre-existing knowledge.

\textbf{Balancing flexibility and regularization in finetuning}. Our investigation reveals a fundamental balance between model flexibility and regularization. Standard finetuning exhibits the highest degree of adaptability but suffers from issues such as low stability and a propensity for model collapse. OFT, despite its simplicity, achieves a compelling compromise between these two aspects by maintaining the pairwise angular relationships among neurons. Additionally, adopting a block-diagonal parameterization acts as an intensified regularizer on the orthogonal matrix.

\textbf{No extra cost during inference}. Distinct from ControlNet, the OFT approach does not introduce any inference-time computational overhead to the finetuned model. During inference, the orthogonal matrix $\bm{R}$ learned during training can be merged with the pretrained weights $\bm{W}^0$, yielding a new equivalent weight matrix $\thickmuskip=2mu \medmuskip=2mu \bm{W}=\bm{R}\bm{W}^0$. Consequently, the inference throughput remains identical to that of the original pretrained model.

\section{Experiments and Results}

\textbf{Experimental setup}. In our empirical studies, Stable Diffusion v1.5~\cite{rombach2022high} serves as the foundational pretrained text-to-image architecture. For an equitable assessment, image samples from each compared method are randomly selected. The protocol for subject-driven image generation largely mirrors that of DreamBooth~\cite{ruiz2023dreambooth}, while the procedures for controllable generation align with ControlNet~\cite{zhang2023adding} and T2I-Adapter~\cite{mou2023t2i}. To guarantee fair evaluation against LoRA, OFT or COFT are confined to the same layers targeted by LoRA. Additional implementation specifics and outcomes are documented in Appendix~\ref{app:settings}.

\subsection{Subject-driven Generation}

\textbf{Experimental details}. DreamBooth~\cite{ruiz2023dreambooth} and LoRA~\cite{hulora2022} constitute the baseline methods in this evaluation. All compared methods employ the loss function adopted in DreamBooth. For both DreamBooth and LoRA, the original methodological guidelines and optimal hyperparameter configurations are followed. Further supplementary results are included in Appendix~\ref{app:settings},\ref{app:coft_oft},\ref{app:more_qualitative},\ref{app:failure}.

\textbf{Finetuning stability and convergence}. To begin, we assess the stability and convergence rate during finetuning for DreamBooth, LoRA, OFT, and COFT. The outcomes, illustrated in Figure~\ref{fig:teaser} and Figure~\ref{fig:db_convergence}, indicate that COFT and OFT consistently maintain stable finetuning behavior when applied to the diffusion model. Notably, after 400 iterations, both DreamBooth and OFT-based methods exhibit effective subject control, whereas LoRA struggles to retain subject identity. When the finetuning process reaches 2000 iterations, DreamBooth begins to produce degenerate images, and LoRA inaccurately synthesizes yellow fur instead of a yellow shirt. On the other hand, OFT and COFT continue to demonstrate robust and consistent control over the output images. These findings underscore the ability of our OFT approach to combine rapid convergence with training stability for subject-driven generation tasks. Importantly, this resilience to the number of finetuning iterations is advantageous as it eliminates the necessity for meticulous tuning of iteration counts across different subjects. Consequently, OFT and COFT permit the use of a relatively high iteration count by default. For COFT, selecting an appropriate $\epsilon$ allows both learning rate and iteration number to be set with minimal effort.

\textbf{Quantitative comparison}. Adopting the evaluation protocol from \cite{ruiz2023dreambooth}, we quantitatively assess subject fidelity (using DINO~\cite{caron2021emerging} and CLIP-I~\cite{radford2021learning}), text prompt fidelity (CLIP-T~\cite{radford2021learning}), and sample diversity (LPIPS~\cite{zhang2018unreasonable}). In this setting, CLIP-I represents the average cosine similarity of CLIP embeddings comparing generated images to real ones, while DINO similarly calculates pairwise similarity using ViT S/16 DINO embeddings. CLIP-T measures the average cosine similarity between the CLIP representations of text prompts and those of generated images. Additionally, we compute the average LPIPS cosine similarity among images generated for the same subject and prompt. According to Table~\ref{table:dreambooth_final}, both COFT and OFT surpass DreamBooth and LoRA by a significant margin in DINO and CLIP-I metrics, and achieve either superior or on-par results in prompt fidelity and diversity. Each method is evaluated by finetuning the same pretrained model across 30 different random seeds to account for variability. These results confirm that our OFT framework not only facilitates improved stability and convergence during training but also consistently leads to higher quality outcomes.

\textbf{Qualitative comparison}. To provide clearer insight into the advantages of OFT, we present several randomly selected samples demonstrating subject-driven generation in Figure~\ref{fig:dreambooth_final}. In order to ensure a fair evaluation, each method generates results from the identical finetuned model, and no individual text prompt optimization is performed for the final outputs. For every method under consideration, the model yielding the highest validation CLIP scores is chosen. As depicted in Figure~\ref{fig:dreambooth_final}, both OFT and COFT consistently exhibit strong semantic preservation of the subject, while LoRA frequently struggles to retain subject identity (\emph{e.g.}, LoRA completely fails to maintain subject characteristics in the bowl instance). Additionally, both OFT and COFT demonstrate more precise text-prompt control compared to DreamBooth, which, although effective at retaining the subject, often does not produce outputs that align well with the textual descriptions (\emph{e.g.}, see the top row of the bowl example). These qualitative observations indicate that OFT can successfully balance both subject fidelity and controllability. Furthermore, OFT demonstrates robustness with respect to the number of iterations, consistently achieving strong performance even with many iterations—an aspect where DreamBooth and LoRA falter. Additional qualitative results are provided in Appendix~\ref{app:more_qualitative}, and human evaluation results presented in Appendix~\ref{app:human} further substantiate the superior performance of OFT.

\subsection{Controllable Generation}

\textbf{Settings}. The baselines utilized for comparison include ControlNet~\cite{zhang2023adding}, T2I-Adapter~\cite{mou2023t2i}, and LoRA~\cite{hulora2022}. Three demanding controllable generation tasks are examined in this study: Canny edge to image~(C2I) on the COCO dataset~\cite{lin2014microsoft}, segmentation map to image~(S2I) on ADE20K~\cite{zhou2017scene}, and landmark to face~(L2F) using CelebA-HQ~\cite{karras2018progressive,xia2021tedigan}. All approaches fine-tune Stable Diffusion~(SD) v1.5 across these datasets for 20 epochs. Further examples and extended analyses are presented in Appendix~\ref{app:more_qualitative},~\ref{app:more_control_task}, and~\ref{app:failure}.

\textbf{Convergence}. We assess how rapidly ControlNet, T2I-Adapter, LoRA, and COFT reach convergence on the L2F task, utilizing both quantitative measurements and visual analysis. For the quantitative metric, we determine the average $\ell_2$ distance between the control face landmarks and those predicted by the models. As illustrated in Figure~\ref{fig:face_convergence}, we present the error in face landmark localization as a function of training epochs for each fine-tuned model. The results indicate that both COFT and OFT converge at a markedly higher rate compared to the alternatives. Notably, LoRA requires 20 epochs to match the level of performance attained by OFT at only the 8-th epoch. It is important to highlight that OFT and COFT contain a comparable number of trainable parameters to LoRA (significantly fewer than ControlNet), yet offer much improved efficiency in terms of convergence speed over prior methods. Additionally, the efficiency of OFT is supported by evidence in Figure~\ref{fig:teaser}; the rightmost example demonstrates OFT’s ability to converge effectively even with just 5\% of the ADE20K dataset, emphasizing its superior data efficiency over ControlNet and LoRA. For the qualitative analysis, we compare OFT, COFT, and ControlNet—since ControlNet exhibits landmark error closest to that of OFT. The visualizations in Figure~\ref{fig:face_convergence_qual} confirm that OFT and COFT both converge reliably, progressively aligning the synthesized face poses with the target landmarks. In contrast, ControlNet shows an abrupt onset of controllability around the 8-th epoch, aligning with the sharp convergence observed in \cite{zhang2023adding}. This consistent and gradual convergence characteristic of OFT facilitates its practical deployment, as its training process exhibits much greater stability and predictability, making the search for effective hyperparameters considerably more straightforward.

\textbf{Quantitative comparison}. To assess the effectiveness of controllable generation, we propose a control consistency metric. This metric involves extracting the control signal from the output image and directly comparing it to the initial input control signal. In the C2I task, we report both IoU and F1 scores; for S2I, we include mean IoU along with mean and overall accuracy; and for L2F, we utilize the mean $\ell_2$ distance calculated between the predicted landmarks and the reference control landmarks. Further details on these consistency measures are provided in Appendix~\ref{app:settings}. For all baselines, optimal hyperparameter settings are adopted to ensure a fair comparison. According to Table~\ref{table:control_gen}, both OFT and COFT outperform all other approaches in terms of control strength and accuracy. We note that adapter-based techniques (\emph{e.g.}, T2I-Adapter and ControlNet) demonstrate slower convergence and ultimately underperform relative to OFT and COFT. LoRA shows superior results compared to ControlNet for the S2I task but lags behind in the C2I and L2F tasks. Notably, even with extensive hyperparameter tuning, LoRA’s performance appears inherently limited. In contrast, OFT exhibits continued improvements as the model size—and, correspondingly, the number of trainable parameters—increases, with no evidence of reaching a performance plateau. Importantly, our quantitative methodology for evaluating controllable generation constitutes a novel aspect of our work, providing a precise assessment of control effectiveness in fine-tuned models, contingent on the performance limits of the segmentation or detection models employed.

\textbf{Qualitative comparison}. We conduct a qualitative analysis of OFT and COFT against leading approaches such as ControlNet, T2I-Adapter, and LoRA. As illustrated by the randomly sampled images in Figure~\ref{fig:control_final}, OFT and COFT consistently produce images that are not only realistic and high in fidelity but also adhere closely to the intended controls. In the S2I task, LoRA is unable to generate images consistent with the provided segmentation map, while ControlNet, OFT, and COFT effectively regulate the output. Compared with ControlNet, OFT and COFT generate images with superior fidelity, richer details, and more plausible geometries, despite utilizing fewer model parameters. For the C2I task, OFT and COFT are able to generate semantically relevant images from approximate Canny edges, a challenge where T2I-Adapter and LoRA underperform. With regard to the L2F task, our framework achieves the most precise pose-controlled facial generation, even under difficult pose conditions. Across all three types of control tasks, our results indicate that OFT and COFT outperform the current state-of-the-art methods in qualitative terms, highlighting the strength of the OFT framework for controllable image synthesis. Further qualitative comparisons covering all three control tasks are available in Appendix~\ref{app:control_exp}, and additional evidence demonstrating OFT’s robust performance on more control scenarios—such as dense pose to human body, sketch to image, and depth to image—can be found in Appendix~\ref{app:more_control_task}.

\section{Concluding Remarks and Open Problems}

Inspired by the crucial role that angular relationships between neurons play in defining visual semantics, this work introduces orthogonal finetuning—a straightforward yet powerful approach for enhancing controllability in text-to-image diffusion models. This method is specifically tailored for two primary text-to-image generation tasks: subject-driven generation and controllable generation. In comparison to existing approaches, OFT offers superior controllability and more stable finetuning dynamics while requiring fewer trainable parameters. Furthermore, OFT does not incur extra inference computational costs, making it a practical solution for deployment.

The development of OFT also raises several compelling research questions. First, the orthogonality constraint in OFT is enforced through Cayley parametrization, which necessitates a matrix inversion and can hinder the scalability of the method. While block diagonal parametrization helps alleviate this limitation, developing efficient, differentiable methods for matrix inversion remains an open challenge. Second, OFT is uniquely poised for compositional applications: orthogonal matrices from multiple OFT-based finetuning instances can be combined via multiplication and the result remains orthogonal. Investigating whether such composed orthogonal matrices successfully retain knowledge from all constituent downstream tasks is a promising research avenue. Lastly, OFT's parameter efficiency is closely tied to the block diagonal architecture, which, while effective, imposes structural bias and restricts expressiveness. Identifying more robust and less biased strategies for enhancing parameter efficiency presents an important unresolved question.